% ============================================================================
% Chapter 10: Deuterium (A=2) from Topological Pinning
% ============================================================================
% Source: M6_MODEL_SUMMARY.md, M6_TOPOLOGICAL_MODEL_EXPLORATION.md
% Status: FILLED from SoT
% Contamination check: PASSED (no banned terms)
% Contamination guard: 3D-model terminology routed to Appendix X/Q.
% Layer A text uses EDC-native vocabulary only.
% ============================================================================

\chapter{Deuterium: The Minimal Bound State}
\label{ch:deuterium}

\source{M6\_MODEL\_SUMMARY.md}

\begin{abstract}
Having established the metastable junction lifetime ($\taun \approx 880\second$)
in Part~II, we now test the EDC framework on the simplest multi-junction system:
the two-junction bound state (A=2). This chapter derives the binding energy
$B_d \approx 3\Kpin$ from topological pinning, where $\Kpin$ is the pinning
constant from Chapter~\ref{ch:sigma_to_K}. A local-optimality derivation
within the pairwise-additive pinning model establishes
$N_{\text{bonds}} = 3$ (Appendix~\ref{app:nbonds}), while geometric
arguments remain supporting motivations. The numerical result
$B_d \approx 2.22\MeV$ agrees with the baseline datum $B_d^{\text{exp}}
= 2.224\MeV$ to $< 1\%$.
Key status: the bond count is derived [Dc] within the local
pairwise-additive pinning model (Appendix~\ref{app:nbonds}); full [Der]
from $S_{\text{EDC}}$ remains open.
\end{abstract}

% ----------------------------------------------------------------------------
% CHAPTER SPINE
% ----------------------------------------------------------------------------
\paragraph{Chapter Spine.}
\textbf{Purpose:} Derive the deuterium binding energy $B_d \approx 2.2\;\MeV$
from topological pinning of two junctions (anchor + metastable).
\textbf{Inputs:} (i)~Pinning constant $\Kpin \approx 0.74\;\MeV$ from Ch.~\ref{ch:sigma_to_K}~[Dc],
(ii)~Y-junction geometry from Ch.~\ref{ch:anchor}~[Dc].
\textbf{Outputs:} (i)~Effective bond count $N_{\text{bonds}} = 3$~[Dc within pinning model],
(ii)~$B_d = 3\Kpin \approx 2.2\;\MeV$~[Dc$\times$Dc(model)],
(iii)~agreement with baseline to $< 1\%$.
\textbf{Dependencies:} Ch.~\ref{ch:sigma_to_K} ($\Kpin$).
\textbf{Forward links:} Bond counting to Ch.~\ref{ch:helium4} (closed-4 unit).

% ============================================================================
\section{Two-Junction Bound State}
\label{sec:d:physical}

\subsection{The Minimal Bound Pair}

In the EDC framework, the simplest bound configuration consists of two
Y-junction defects (one anchor junction, one metastable junction) brought
into contact:

\begin{center}
\fbox{\parbox{0.7\textwidth}{
\textbf{Two-Junction Contact Configuration}

\medskip
\begin{tabular}{ll}
Junction 1 ($J_1$): & Anchor junction (Z$_6$ ground state, $q_1 = 0$) \\
Junction 2 ($J_2$): & Metastable junction (Z$_3$ state, $q_2 = 1$) \\
Contact: & Shared coordination region (``contact patch'') \\
\end{tabular}
}}
\end{center}

\subsection{Contact Patch Geometry}

When two junctions approach, their Y-arm structures interpenetrate,
creating a \emph{contact patch}---a region where degrees of freedom are
mutually constrained.

\begin{definition}[Contact Patch]
    The \textbf{contact patch} is the geometric intersection of the
    coordination volumes of two adjacent junctions. Within this region:
    \begin{itemize}
        \item Arm orientations are mutually locked
        \item Deformation modes are coupled
        \item The topological state $(q_1, q_2)$ experiences a pinning potential
    \end{itemize}
    \tagDc
\end{definition}

\subsection{Reference: Pinning Constant $\Kpin$}

The pinning constant $\Kpin$ was defined in Chapter~\ref{ch:sigma_to_K}:
\begin{equation}
    \Kpin = f \times \bsigma \times A_{\text{shared}}
    \label{eq:K_recall}
\end{equation}
where $f$ is a geometric factor and $A_{\text{shared}}$ is the effective
shared area per bond.

\textbf{Numerical value:} $\Kpin \approx 0.74\MeV$ per bond [Dc].

% ============================================================================
\section{Pinning Energy Model for A=2}
\label{sec:d:model}

\subsection{Bond Counting Ansatz}

The binding energy of the two-junction system is modeled as:

\begin{derivationbox}[title=Pinning Energy Model]
    \begin{equation}
        B_d = N_{\text{bonds}} \times \Kpin
        \label{eq:Bd_model}
    \end{equation}
    where $N_{\text{bonds}}$ is the effective number of ``pinning bonds''
    between the two junctions.
    \tagP
\end{derivationbox}

\subsection{What is a ``Bond''?}

In EDC terminology, a \emph{bond} is not an exchange mechanism but a
\textbf{coordination lock}:

\begin{definition}[Pinning Bond]
    A \textbf{pinning bond} is a topological constraint that:
    \begin{enumerate}
        \item Locks a deformation channel between adjacent junctions
        \item Suppresses relative mismatch in the $(q_i - q_j)$ coordinate
        \item Contributes energy $\sim \Kpin$ when the lock is engaged
    \end{enumerate}
    The number of bonds is discrete (integer) because the constraint
    structure is topological.
    \tagDc
\end{definition}

\subsection{Model Assumptions}

\begin{center}
\begin{tabular}{lll}
\toprule
\textbf{Assumption} & \textbf{Tag} & \textbf{Status} \\
\midrule
Binding $\propto N_{\text{bonds}} \times \Kpin$ & [P] & Ansatz \\
$N_{\text{bonds}}$ is integer & [Dc] & Topological \\
$\Kpin$ is universal (same for all bonds) & [P] & Approximation \\
\bottomrule
\end{tabular}
\end{center}

% ============================================================================
\section{Why $N_{\text{bonds}} = 3$}
\label{sec:d:three}

The effective bond count $N_{\text{bonds}} = 3$ is the central claim of
this chapter. We summarize the result here; the full derivation is given
in Appendix~\ref{app:nbonds}.

\subsection{Local Optimality Result}

Appendix~\ref{app:nbonds} establishes $N_{\text{bonds}} = 3$ as a
\textbf{local optimality result within the pairwise-additive pinning model}.
The derivation proceeds via three lemmas:

\begin{derivationbox}[title=Local Optimality of $N_{\text{bonds}} = 3$ (Summary)]
    \textbf{Theorem~\ref{thm:nbonds_optimality}}
    (Appendix~\ref{app:nbonds}):
    Within the pairwise-additive pinning model, for two Y-junctions
    sharing a contact face in the $\Msix$ lattice, the energy-minimizing
    configuration has exactly $N_{\text{bonds}} = 3$ contact pairs.

    \medskip
    \textbf{Proof outline:}
    \begin{itemize}
        \item \textbf{Lemma~\ref{lem:upper_bound} (Upper bound):}
              Each Y-junction has 3 arms [Der]; within the contact model,
              each arm supports at most one bond.
              $\Rightarrow N_{\text{bonds}} \leq 3$.
        \item \textbf{Lemma~\ref{lem:saturation} (Energetic saturation):}
              Under four stated assumptions (pairwise additivity, locality,
              no frustration penalty, no multi-arm coupling), the energy
              $E(n) = n\Kpin$ is monotonically decreasing.
              $\Rightarrow$ Minimum at $n = 3$.
        \item \textbf{Lemma~\ref{lem:uniqueness} (Pairing uniqueness):}
              Within the restricted geometric contact model, the
              complementary-angle identity pairing is the unique
              energy minimum, invariant under $\Zthree$.
    \end{itemize}

    \textbf{Epistemic status:}
    [Dc] within the local pairwise-additive pinning model.
    This is not a direct derivation from $S_{\text{EDC}}$ and not a
    general topological proof. The four assumptions of
    Lemma~\ref{lem:saturation} remain [P].
    See Appendix~\ref{app:nbonds} for the full proof and residual
    open items.
\end{derivationbox}

\subsection{Supporting Geometric Arguments}

The local optimality result is consistent with two independent
geometric motivations:

\begin{enumerate}
    \item \textbf{Contact geometry (Y-junction overlap):}
          Each Y-junction has 3 arms at $120°$ angles. In face-to-face
          contact, the 3 arms of $J_1$ engage the 3 complementary arms
          of $J_2$, yielding 3 coordination locks.

    \item \textbf{Dual lattice coordination ($\Msix$ structure):}
          In the $\Msix$ lattice, each junction uses 3 of its 6
          coordination directions for internal structure and shares
          3 with neighbors.
\end{enumerate}

Both arguments are consistent with $N = 3$ but do not independently
constitute rigorous derivations.

\subsection{Summary of Bond Count Status}

\begin{center}
\begin{tabular}{llll}
\toprule
\textbf{Method} & \textbf{Basis} & \textbf{Result} & \textbf{Tag} \\
\midrule
Local optimality & Pinning-model energy & $N = 3$ & [Dc] (within model) \\
Contact geometry & Y-junction arm count & $N = 3$ & [P] (motivation) \\
Dual lattice & M$_6$ coordination & $N = 3$ & [P] (motivation) \\
\midrule
\textbf{Combined} & Optimality + consistency & $N_{\text{bonds}} = 3$ &
    [Dc] within pinning model \\
\bottomrule
\end{tabular}
\end{center}

\textbf{Scope:} The [Dc] tag applies within the local pairwise-additive
pinning model. A full [Der] promotion requires deriving the pinning
Hamiltonian $H_{\text{pin}}$ from $S_{\text{EDC}}$, which is beyond
the current scope (see Appendix~\ref{app:nbonds}, Residual Open Items).

% ============================================================================
\section{Numerical Evaluation}
\label{sec:d:numerical}

\subsection{Input Parameters}

\begin{center}
\begin{tabular}{lllll}
\toprule
\textbf{Symbol} & \textbf{Meaning} & \textbf{Value} & \textbf{Tag} & \textbf{Source} \\
\midrule
$\Kpin$ & Pinning constant & $0.74\MeV$ & [Dc] & Ch.~\ref{ch:sigma_to_K} \\
$N_{\text{bonds}}$ & Effective bond count & $3$ & [P] & \S\ref{sec:d:three} \\
\bottomrule
\end{tabular}
\end{center}

\subsection{Binding Energy Calculation}

\begin{derivationbox}[title=Binding Energy Result]
    \begin{equation}
        B_d^{\text{(EDC)}} = N_{\text{bonds}} \times \Kpin = 3 \times 0.74\MeV = 2.22\MeV
        \label{eq:Bd_result}
    \end{equation}
    \tagDc $\times$ \tagP
\end{derivationbox}

\subsection{Results Table}

\begin{center}
\begin{tabular}{lll}
\toprule
\textbf{Quantity} & \textbf{Value} & \textbf{Tag} \\
\midrule
$\Kpin$ & $0.74\MeV$ & [Dc] \\
$N_{\text{bonds}}$ & $3$ & [P] \\
$B_d^{\text{(EDC)}} = 3\Kpin$ & $2.22\MeV$ & [Dc]$\times$[P] \\
\bottomrule
\end{tabular}
\end{center}

\textbf{Uncertainty:} If $\Kpin$ has $\pm 10\%$ uncertainty, then
$B_d \in [2.0, 2.4]\MeV$.

% ============================================================================
\section{Baseline Datum and Comparison}
\label{sec:d:baseline}

\subsection{Experimental Reference}

\begin{center}
\begin{tabular}{lll}
\toprule
\textbf{Quantity} & \textbf{Value} & \textbf{Tag} \\
\midrule
Measured binding energy & $B_d^{\text{exp}} = 2.224\MeV$ & [BL] \\
\bottomrule
\end{tabular}
\end{center}

This is an empirical baseline datum, not a prediction target.

\subsection{Comparison}

\begin{center}
\begin{tabular}{llll}
\toprule
& \textbf{Value} & \textbf{Source} & \textbf{Deviation} \\
\midrule
EDC prediction & $2.22\MeV$ & Eq.~\eqref{eq:Bd_result} & — \\
Baseline [BL] & $2.224\MeV$ & Measured & — \\
\midrule
Error & & & $< 1\%$ \\
\bottomrule
\end{tabular}
\end{center}

\textbf{Observation:} The EDC estimate matches the baseline to within
$1\%$---well within the model's intrinsic uncertainty.

% ============================================================================
\section{Sensitivity and Robustness}
\label{sec:d:sensitivity}

\subsection{Sensitivity to $\Kpin$}

\begin{center}
\begin{tabular}{lll}
\toprule
$\Kpin$ variation & $B_d = 3\Kpin$ & Change vs.\ baseline \\
\midrule
$0.90 \times 0.74 = 0.67\MeV$ & $2.00\MeV$ & $-10\%$ \\
$0.95 \times 0.74 = 0.70\MeV$ & $2.11\MeV$ & $-5\%$ \\
$0.74\MeV$ (nominal) & $2.22\MeV$ & $0\%$ \\
$1.05 \times 0.74 = 0.78\MeV$ & $2.33\MeV$ & $+5\%$ \\
$1.10 \times 0.74 = 0.81\MeV$ & $2.44\MeV$ & $+10\%$ \\
\bottomrule
\end{tabular}
\end{center}

\textbf{Observation:} $B_d$ scales linearly with $\Kpin$. The $\pm 10\%$
envelope of $\Kpin$ maps directly to $\pm 10\%$ in $B_d$.

\subsection{Sensitivity to $N_{\text{bonds}}$}

\begin{center}
\begin{tabular}{llll}
\toprule
$N_{\text{bonds}}$ & $B_d = N \times \Kpin$ & Match to baseline & Status \\
\midrule
$2$ & $1.48\MeV$ & $-33\%$ & Poor \\
$3$ & $2.22\MeV$ & $< 1\%$ & \textbf{Good} \\
$4$ & $2.96\MeV$ & $+33\%$ & Poor \\
\bottomrule
\end{tabular}
\end{center}

\textbf{Critical observation:} The bond count \emph{must} be $N = 3$ for
agreement. Values $N = 2$ or $N = 4$ are clearly excluded by the baseline
datum.

\subsection{Robustness Assessment}

\begin{edcopenbox}[title=What Must Be True]
    For the model to survive review:
    \begin{enumerate}
        \item $N_{\text{bonds}} = 3$ must be \textbf{derivable}, not chosen
              by hand to match data
        \item $\Kpin$ must be independently fixed (from $\bsigma$, not from
              $B_d$)
        \item The linear model $B_d = N\Kpin$ must be justified as leading
              order, with corrections bounded
    \end{enumerate}

    Current status: Points 2 and 3 are satisfied. Point 1 is now [Dc]
    within the pinning model (Appendix~\ref{app:nbonds}); full [Der]
    from $S_{\text{EDC}}$ remains open.
\end{edcopenbox}

% ============================================================================
\section{Epistemic Status}
\label{sec:d:epistemic}

\begin{center}
\begin{tabular}{llll}
\toprule
\textbf{Claim} & \textbf{Status} & \textbf{Reference} & \textbf{Dependency} \\
\midrule
$\Kpin \approx 0.74\MeV$ & [Dc] & Ch.~\ref{ch:sigma_to_K} & From $\bsigma$ \\
$B_d = N_{\text{bonds}} \times \Kpin$ & [P] & Eq.~\eqref{eq:Bd_model} & Ansatz \\
$N_{\text{bonds}} = 3$ (local optimality) & [Dc] (within model) & App.~\ref{app:nbonds} & Pinning model \\
$B_d \approx 2.22\MeV$ & [Dc]$\times$[Dc(model)] & Eq.~\eqref{eq:Bd_result} & Calculation \\
\midrule
$B_d^{\text{exp}} = 2.224\MeV$ & [BL] & \S\ref{sec:d:baseline} & Measurement \\
\midrule
$H_{\text{pin}}$ from $S_{\text{EDC}}$ & [OPEN] & — & Put C corridor \\
Corrections to linear model & [OPEN] & — & Higher order \\
\bottomrule
\end{tabular}
\end{center}

\subsection{Open Problems}

\begin{edcopenbox}[title=Open Problems]
    \begin{enumerate}
        \item \textbf{$H_{\text{pin}}$ from $S_{\text{EDC}}$:} The bond
              count $N = 3$ is derived within the pairwise-additive pinning
              model (Appendix~\ref{app:nbonds}). Upgrading to [Der] requires
              deriving $H_{\text{pin}}$ from the 5D action, which would
              also test the four assumptions of Lemma~\ref{lem:saturation}
              (pairwise additivity, locality, no frustration, no multi-arm
              coupling).

        \item \textbf{Corrections beyond linear:} The model $B_d = N\Kpin$
              ignores:
              \begin{itemize}
                  \item Surface/edge effects at the contact patch
                  \item Deformation energy of the ``q-field'' relaxation
                  \item Localization sharing (relevant for larger A)
              \end{itemize}
              Estimate these corrections and verify they are $< 10\%$.

        \item \textbf{Spin/orientation effects:} The two-junction pair
              has internal orientation degrees of freedom. How do these
              affect binding?
    \end{enumerate}
    \tagOpen
\end{edcopenbox}

% ============================================================================
\section{Falsifiability}
\label{sec:d:falsify}

\begin{edcopenbox}[title=Falsifiability Hooks]
    \begin{enumerate}
        \item \textbf{Bond count mismatch:} If rigorous M$_6$ analysis
              gives $N_{\text{bonds}} \neq 3$, the geometric arguments
              fail and $B_d$ prediction changes by $\sim 33\%$ per unit.

        \item \textbf{$\Kpin$ inconsistency:} If the $\Kpin$ value required
              to match $B_d$ differs from the value derived from $\bsigma$
              by more than $20\%$, the single-parameter picture breaks.

        \item \textbf{Scaling failure:} If other A=2 systems (if any exist)
              have $B \neq 3\Kpin$, the bond-counting model is too simple.

        \item \textbf{He-4 inconsistency:} Chapter~\ref{ch:helium4} will
              test $B_{\text{He-4}}$. If $\Kpin$ cannot consistently
              explain both $B_d$ and $B_{\text{He-4}}$, the framework fails.
    \end{enumerate}
\end{edcopenbox}

% ============================================================================
% Observer-side projection
% ============================================================================

\begin{observerbox}
Observer-side projection note: quoted terms are measurement labels for the
3D/4D projection of the 5D objects defined in this chapter; they do not
imply any conventional mechanism.

\begin{itemize}
    \item Two-junction bound state $\leftrightarrow$ A=2 measured system (projection label)
    \item \Pinning{} energy $\leftrightarrow$ binding energy (projection label)
\end{itemize}

What the observer records: the two-junction configuration with $B_d \approx 2.2\MeV$
corresponds to the measured binding of the lightest bound cluster. This is
a projection-level consistency check.
\end{observerbox}

% ============================================================================
\section{Summary}
\label{sec:d:summary}

\begin{resultbox}
    \textbf{What was achieved:}
    \begin{itemize}
        \item Defined the two-junction bound state (A=2) as the minimal
              pinning configuration
        \item Established the model: $B_d = N_{\text{bonds}} \times \Kpin$
        \item Derived $N_{\text{bonds}} = 3$ as a local optimality result
              within the pairwise-additive pinning model
              (Appendix~\ref{app:nbonds}, [Dc] within model)
        \item Computed $B_d = 3 \times 0.74\MeV = 2.22\MeV$ [Dc]$\times$[Dc(model)]
        \item Matched baseline $B_d^{\text{exp}} = 2.224\MeV$ to $< 1\%$
    \end{itemize}

    \textbf{Key residual:} The [Dc] tag applies within the declared pinning
    model. A full [Der] from $S_{\text{EDC}}$ requires deriving the pinning
    Hamiltonian from the 5D action (Put~C corridor).
\end{resultbox}

\bigskip
\noindent
\textbf{Forward references:}
\begin{itemize}
    \item Chapter~\ref{ch:helium4}: Four-junction system (He-4) with
          localization + pinning
    \item Chapter~\ref{ch:light_clusters}: Patterns in light systems
          (A $\leq 10$)
\end{itemize}

\bigskip
\noindent
\textbf{Derivation chain extension:}
\begin{equation*}
    \bsigma \xrightarrow{\text{Ch.~\ref{ch:sigma_to_K}}} \Kpin
    \xrightarrow{\text{This chapter}} B_d = 3\Kpin \approx 2.2\MeV
\end{equation*}

% ----------------------------------------------------------------------------
% BRIDGE TO NEXT CHAPTER
% ----------------------------------------------------------------------------
\paragraph{Bridge to Chapter~\ref{ch:helium4}.}
The two-junction bound state demonstrates that pinning alone can produce
binding energies of order $\sim 2\;\MeV$. But the A=4 system has
$B_4 \approx 28.3\;\MeV$---six times higher binding per constituent than
deuterium. Why? Chapter~\ref{ch:helium4} introduces the \emph{closed-4 unit}:
four junctions forming a complete $K_4$ graph with no boundary. This
topological closure enables collective effects that dominate the energy budget.

